\appendix

\section{Interpretation of $(I + \alpha L)^{-1}$}\label{sec:operator_interpretation}
Throughout the paper we have used $K^{-1} = (I + \alpha L)^{-1}$ as our
long-range propgation operator. We now show mathatically that $K^{-1}$ indeed
does multi-hop propgation of the normalized adjecacny matrix $\hat A$.

Notice that $K = I + \alpha L = (1 + \alpha)I - \alpha \hat A$ and so using Neumann expansion
\begin{equation}
    K^{-1} = \frac{1}{1 + \alpha}\left(I - \frac{\alpha}{1 + \alpha} \hat A \right)^{-1}=  \frac{1}{1 + \alpha}\sum_{k = 0}^{\infty}\left(\frac{\alpha}{1 + \alpha}\right)^{k} \hat A^k
\end{equation}

The usage of Neumann expansion is valid since for $\alpha > 0$
\begin{equation}
    \rho\left(\frac{\alpha}{1 + \alpha} \hat A \right) = \frac{\alpha}{1 + \alpha} \rho\left( \hat A \right) \leq  \frac{\alpha}{1 + \alpha} < 1
\end{equation}

A conventional GCN layer propgates using $\hat A$
\begin{equation*}
    H^{(l + 1)} = \sigma(\hat A H^{(l)} W^{(l)})
\end{equation*}

This only does 1-hop propagation.

Replacing $\hat A$ with $K^{-1}$ is equivilant to

\begin{equation*}
    H^{(l + 1)} = \sigma(K^{-1} H^{(l)} W^{(l)}) = \sigma\left(\frac{1}{1 + \alpha}\sum_{k = 0}^{\infty}\left(\frac{\alpha}{1 + \alpha}\right)^{k} \hat A^k H^{(l)} W^{(l)}\right)
\end{equation*}

Define
\begin{equation}
    w_k = \frac{1}{1 + \alpha} \left( \frac{\alpha}{1 + \alpha}\right )^k
\end{equation}

Notice that larger $\alpha$ corresponds to slower weight decay. Therefore we can
achive longer effictive propagation by increasing $\alpha$.

\section{Additional Method Details}
\subsection{PCG Algorithm}
\begin{algorithm}[H]
    \capation{Preconditioned Conjugate Gradient}
    \label{alg:pcg}
    \begin{algorithmic}[1]
        \Require{SPD Matrix $A$, RHS $b$, preconditioner $P$, initial guess $x_0$, tolerance $\epsilon$, maximum iterations $K$}

        \State $x \gets x_0$
        \State $r \gets b - Ax$
        \State $z \gets P(r)$
        \State $p \gets z$

        \For{$k = 0, \dots, K - 1$}
        \State $q \gets A p$

        \State $r' \gets \langle r,z \rangle$

        \State $\alpha \gets \frac{r'}{\langle p,q\rangle}$

        \State $x \gets x + \alpha p$
        \State $r \gets r - \alpha q$

        \If{$\frac{\lVert r\rVert_2}{\lVert b\rVert_2} < \epsilon $}
        \State \Return $x$
        \EndIf

        \State $z \gets P(r)$
        \State $\beta \gets \frac{\langle r, z\rangle}{r'}$

        \State $p \gets z + \beta p$

        \EndFor

        \State \Return $x$
    \end{algorithmic}
\end{algorithm}

\subsection{\texttt{MPSSM-PCG} implementation}
% def forward(self, U,S):
%     # X = torch.zeros(U.shape[0], self.W_supp.out_features, dtype= U.dtype, device = U.device)
%     X = torch.zeros(U.shape[0], self.W.out_features, dtype= U.dtype, device = U.device)
%
%     UB = self.B(U)
%
%
%     for _ in range(self.num_steps):
%         # X = self.W_supp(S_supp @ X) + self.W_val(S_val @ X) + UB
%          X =  self.W(S @ X) + UB
%
%     return self.mlp(X)

\begin{algorithm}[H]
    \capation{MPSSM Block}
    \label{alg:mpssm_block}
    \begin{algorithmic}[1]
        \Require{Propgation matrix $S$, Matrix Features $U$, SSM steps $T$}

        \State $X_0 \gets \mathbf{0}$
        \For{$s = 0,\dots T -1$}
        \State $X_{s + 1} \gets SX_{s}W + UB$
        \EndFor

        \State \Return $\texttt{MLP}(X_T)$
    \end{algorithmic}
\end{algorithm}

% H = self.norms[0](self.in_block(U,S))
%
%
%
% for i, block in enumerate(self.blocks):
%     H_t = block(H,S) 
%     H_t = F.dropout(H_t, p = self.dropout, training=self.training)
%     H = self.norms[i + 1](H + H_t)
%
% return H

\begin{algorithm}[H]
    \capation{MPSSM}
    \label{alg:mpssm}
    \begin{algorithmic}[1]
        \Require{Propgation matrix $S$, Matrix Features $U$, SSM steps $T$, SSM Blocks $B$, Dropout $p$}
        \State $H_0 \gets \texttt{GraphNorm}(\texttt{MPSSM}_0(U,S))$

        \For{$b = 0, \dots ,B - 1$}
        \State $H_{b + 1} \gets \texttt{GraphNorm}(H_b + \texttt{Dropout}_p(\texttt{MPSSM}_b(H_b,S)))$
        \EndFor

        \State \Return $H_{B}$

    \end{algorithmic}
\end{algorithm}

% c = 0.1
%
% S= graph_shift(K)
% # S_supp = graph_shift(K)
% # S_val = weighted_graph_shift(K)
%
% features, scale = matrix_node_features(K)
%
% H = self.mpssm(S, features)
%
% src, dst = lower_edge_pattern(K)
% L = torch.zeros_like(K)
%
% if src.numel() > 0:
%     edge_val = (K[src, dst] / scale).unsqueeze(-1)
%
%     edge_features = torch.cat(
%             [
%                 H[src],
%                 H[dst],
%                 edge_val
%             ],
%             dim = -1
%     )
%
%     values = self.factor_decoder(edge_features).squeeze(-1)
%
%     L[src,dst] = values
%
% diag_delta = self.diag_decoder(H).squeeze(-1)
% L_diag =  torch.exp(c * diag_delta)
%
% L = L + torch.diag(L_diag)

\begin{algorithm}[H]
    \capation{MPSSM-PCG}
    \label{alg:mpssm_pcg}
    \begin{algorithmic}[1]
        \Require{System $K$}
        \State $S \gets D^{-1/2}A_SD^{-1/2}$ \Comment{$A_s$ is the sparisty pattern (Eq.\ref{eq:sparsity_pattern}) of $K$ }
        \State $U \gets \texttt{matrix\_node\_features}(K)$ \Comment{$U = [u_1,\dots,u_n]^\top$ when $u_i$ is defined as in Eq.\ref{eq:system_feature_vec}}
        \State $H \gets \texttt{MPSSM}(S,U)$
        \State $E_0 \gets \{(i,j) | i > j, K_{ij} \neq 0\}$
        \State $L_\theta \gets \mathbf{0}$
        \ForAll{$(i,j) \in E_{0}$}
        \State $e_{ij}\gets [H_i || H_j || K_{ij} / s]$
        \State $(L_\theta)_{ij} = \texttt{MLP}_{\text{edge}}(e_{ij})$
        \EndFor

        \For{$i = 1,\dots,n$}
        \State $\delta_i \gets \texttt{MLP}_{\text{diag}}(H_i)$
        \State $(L_\theta)_{ii} \gets e^{c\delta_i}$ \Comment{In our implementation $c = 0.1$}
        \EndFor

        \State \Return $L_\theta$
    \end{algorithmic}
\end{algorithm}

% def forward(self, X : torch.Tensor, A_norm : torch.Tensor = None, preconditioner : Preconditioner = None):
%     if A_norm is None and preconditioner is None:
%         raise ValueError(f"You need a preconditioner or a shift operator")
%
%     H = self.encoder(X) 
%
%
%     for layer,norm in zip(self.layers, self.norms):
%
%         residual = H
%
%         if A_norm is not None:
%             H = A_norm @ H
%         else:
%             H = preconditioner.apply(H)  
%
%         H = layer(H)
%         H = norm(H)
%         H = F.relu(H)
%
%         H = F.dropout(H, p = self.dropout,training=self.training)
%
%         H = H + residual
%
%
%     return self.out(H)
%

% def apply(self, r):
%     is_vec = r.dim() == 1
%
%     if is_vec:
%         r = r.unsqueeze(-1)
%
%
%     y = torch.linalg.solve_triangular(self.L , r, upper = False)
%     y = torch.linalg.solve_triangular(self.L.T , y, upper = True)
%
%     if is_vec:
%         y = y.squeeze(-1)
%
%
%     return y

\begin{algorithm}[H]
    \capation{Apply MPSSM-PCG Preconditioner}
    \label{alg:mpssm_pcg_preconditioner}
    \begin{algorithmic}[1]
        \Require{Sparse Factor $L$, Residual $r$}

        \State $y \gets \texttt{solve\_triangular}(L, r, \text{upper} = F)$ \Comment{Pytorch's \texttt{torch.linalg.solve\_triangular}}
        \State $z \gets \texttt{solve\_triangular}(L^\top, y, \text{upper} = T)$
        \State \Return $z$
    \end{algorithmic}
\end{algorithm}

\subsection{GCN Architicture}
\begin{algorithm}[H]
    \capation{GCN For ECHO-SSSP}
    \label{alg:gcn}
    \begin{algorithmic}[1]
        \Require{Normalized adjecacny $\hat A$, Node Embeddings $X$, Layers $L$, Dropout $p$}
        $H_0 \gets \texttt{MLP}_{\text{encoder}}(X)$
        \For{$i = 0,\dots,L - 1$}
        \State $R_t\gets H_t$
        \State $H_t \gets \texttt{ReLU}(\texttt{LayerNorm}(\hat AH_tW_t))$
        \State $H_t \gets \texttt{Dropout}_p(H_{t}) + R_t$
        \EndFor

        \State \Return $\texttt{MLP}_{\text{out}}(H_L)$
    \end{algorithmic}
\end{algorithm}

\begin{algorithm}[H]
    \capation{PGCN For ECHO-SSSP}
    \label{alg:pgcn}
    \begin{algorithmic}[1]
        \Require{Sparse Factor $L_\theta$, Node Embeddings $X$, Layers $L$, Dropout $p$}
        $H_0 \gets \texttt{MLP}_{\text{encoder}}(X)$
        \For{$i = 0,\dots,L - 1$}
        \State $R_t\gets H_t$
        \State $H_t \gets \texttt{ReLU}(\texttt{LayerNorm}(\texttt{apply}(L_\theta,H_t)W_t))$ \Comment{When \texttt{apply} as in alg.\ref{alg:mpssm_pcg_preconditioner}}
        \State $H_t \gets \texttt{Dropout}_p(H_{t}) + R_t$
        \EndFor

        \State \Return $\texttt{MLP}_{\text{out}}(H_L)$
    \end{algorithmic}
\end{algorithm}

\section{Evaluation}
\subsection{Hardware Configuration}
We run all the tests on Ryzen 9600X CPU, 32GB of RAM and an AMD Radeon 9070 XT GPU with 16 GB of VRAM.

\subsection{Software and Libraries}
Specific libraries are provided using `requirements.txt`. 

\subsection{Expirements Hyperparameters}
\begin{table}[H]
    \centering
    \label{tbl:hyperparams_mpssmpcg}
    \caption{Hyperparameters used in MPSSM-PCG ECHO-SSSP/Peptides-struct preconditioning expirement}
    \small
    \begin{tabular}{c c}
        Hyperparameter       & Value                \\
        \hline
        Optimizer            & Adam                 \\
        Learning Rate        & $10^{-3}$            \\
        Train Epochs         & 50                   \\
        Checkpoint Selection & Best Validation loss \\
        Hidden Dimesntion    & 16                   \\
        Blocks               & 2                    \\
        Steps                & 4                    \\
        Normalization        & \texttt{GraphNorm}   \\
        Dropout              & 0                    \\
    \end{tabular}
\end{table}

\begin{table}[H]
    \centering
    \caption{NeuralIF Hyperparameters ECHO-SSSP/Peptides-struct preconditioning expirement}
    \small
    \begin{tabular}{c c}
        Hyperparameter              & Value                \\
        \hline
        Optimizer                   & Adam                 \\
        Learning Rate               & $10^{-3}$            \\
        Train Epochs                & 50                   \\
        Checkpoint Selection        & Best Validation loss \\
        message\_passing\_steps     & 4                    \\
        global\_features Dimesntion & 0                    \\
        latent\_size                & 8                    \\
        skip\_connections           & True                 \\
        activation                  & ReLU                 \\
        aggregate                   & None                 \\
        decode\_nodes               & False                \\
        normalize\_diag             & False                \\
        graph\_norm                 & True                 \\
        two\_hop                    & False                \\
        edge\_features              & 1
    \end{tabular}
\end{table}

This configuration is the closest configuation to what the authors of NeuralIF
\cite{hausnerNeuralIncompleteFactorization2024} chose. We just used 4 message
passing steps instead of 3 to have a close parameter count to ours.

\begin{table}[H]
    \centering
    \caption{Hyperparameters used in MPSSM-PCG ECHO-SSSP expirement}
    \small
    \begin{tabular}{c c}
        Hyperparameter       & Value                \\
        \hline
        Optimizer            & Adam                 \\
        Learning Rate        & $10^{-4}$            \\
        Train Epochs         & 30                   \\
        Checkpoint Selection & Best Validation loss \\
        Hidden Dimesntion    & 16                   \\
        Blocks               & 2                    \\
        Steps                & 4                    \\
        Normalization        & \texttt{GraphNorm}   \\
        Dropout              & 0                    \\
    \end{tabular}
\end{table}

\begin{table}[H]
    \centering
    \caption{GCN hyperparameters for ECHO-SSSP long-range propagation expirements}
    \small
    \begin{tabular}{c c}
        Hyperparameter       & Value                \\
        \hline
        Optimizer            & Adam                 \\
        Learning Rate        & $10^{-5}$            \\
        Train Epochs         & 300                  \\
        Checkpoint Selection & Best Validation loss \\
        Hidden Dimenstion    & 384                  \\
        Normalization        & \texttt{LayerNorm}   \\
        Dropout              & 0.1                  \\
    \end{tabular}
\end{table}

GCN layers used in the expirements are stated in each table.

The total parameter count of $\texttt{MPSSM-PCG}$ is 2594 while $\texttt{NeuralIF}$ is 2488.

\subsection{Metrices}
\paragraph{Setup time}\label{par:setup_time} Measures the cost
of contructing the preconditioner for a unseen system $K$. For
$\texttt{MPSSM-PCG}$, this consists of constructing the matrix representation,
running the MPSSM encoder and decoding $L_\theta$. Simliarly, for
\texttt{NeuralIF} this is the forward pass for constructing $L_\theta$.
Preconditiong training is not part of the setup time.

Since the tested preconditioners only rely on $K$, for each graph we sample $N_{rhs} = 5$ random RHS and reuse the computed
preconditioner to solve the 5 systems, so the actual setup time presented in tables is
\begin{equation}
    t_{setup} = \frac{t_{\text{preconditioner construction}}}{N_{rhs}}
\end{equation}

\paragraph{Solve Time}\label{par:solve_time} Measures the execution time of PCG
\ref{alg:pcg} and the applications of the preconditioner. MPSSM-PCG and NeuralIF
preconditioners are applied using \ref{alg:mpssm_pcg_preconditioner}.

\paragraph{Total Time}\label{par:total_time} The sum of the setup time and solve time.

\paragraph{Mean Squared Error} For a
graph $G$ with predictions $\mathbf{\hat y}$  and targets $\mathbf{y}$ containing $N$ targets
\begin{equation}\label{eq:mse}
    \text{MSE}(G) =  \frac{1}{N}\sum_{i = 1}^{N} \ (y_i - \hat y_i)^2
\end{equation}

for a mini batch $B$

\begin{equation}
    \text{MSE}(B) = \frac{1}{|B|}\sum_{G \in B} \text{MSE}(G)
\end{equation}

\paragraph{Mean Absolute Error}
\begin{equation}\label{eq:mae}
    \text{MAE} = \frac{\sum_G \sum_{v\in V_G} |y_v - \hat y_v|}{\sum_{G} |V_G|}
\end{equation}

During training, the training loss gives equal weights to each graph
(regardelss of the number of nodes). In evaluation MAE give equal weight to
each target node.

\subsection{Additional Expirements}
In addition to Tabel \ref{tbl:echo_sssp_preconditioning_time} we include both the relative residual
\begin{equation}
    \frac{
        \lVert b - K\hat x\rVert_2
    } {
        \lVert b\rVert_2
    }
\end{equation}

and the relative solution error

\begin{equation}
    \frac{
        \lVert x - \hat x\rVert_2
    } {
        \lVert x\rVert_2
    }
\end{equation}

\begin{table}[H]
    \centering
    \caption{
        Preconditioning performence on ECHO-SSSP graph systems. Mehods are
        evaluated on their PCG performence on unseen test system. Results are
        averaged over three seeds $\{0,1,2\}$.
    }
    \small
    \label{tbl:echo_sssp_preconditioning_err}
    \begin{tabular}{c c c}
        Method           & Rel. Error  $\downarrow$                  & Rel. Reisdual  $\downarrow$               \\
        \hline
        Identity         & $2.094\cdot10^{-6}\pm2.305\cdot{10^{-9}}$ & $7.410\cdot10^{-7}\pm1.065\cdot{10^{-9}}$ \\
        Jacobi           & $2.112\cdot10^{-6}\pm7.335\cdot{10^{-9}}$ & $7.348\cdot10^{-7}\pm1.985\cdot{10^{-9}}$ \\
        NeuralIF         & $1.286\cdot10^{-6}\pm6.873\cdot{10^{-9}}$ & $4.725\cdot10^{-7}\pm1.534\cdot{10^{-9}}$ \\
        MPSSM-PCG (ours) & $1.289\cdot10^{-6}\pm1.076\cdot{10^{-8}}$ & $4.725\cdot10^{-7}\pm2.207\cdot{10^{-9}}$ \\
    \end{tabular}
\end{table}

\begin{table}[H]
    \caption{
        Preconditioning performence on peptides-struct graph systems $K = I +
            8L$, with 2000/500/500 split. Mehods are evaluated on their PCG
        performence on unseen test system. Results are averaged over three
        seeds $\{0,1,2\}$. Timing are presented in ms.
    }
    \centering
    \small
    \label{tbl:peptides_struct_preconditioning_time}
    \begin{tabular}{c c c c c c}
        Method   & Conv. \%   $\uparrow$ & PCG Iteration $\downarrow$ & Setup Time (ms) $\downarrow$ & Solve Time (ms)$\downarrow$ & Total Time (ms)  $\downarrow$ \\
        \hline
        Identity & $100.0\pm0.0$         & $21.90\pm0.00$             & $0.001\pm0.000 $         & $2.244\pm0.002 $        & $2.245\pm0.002$           \\
        Jacobi   & $100.0\pm0.0$         & $21.64\pm0.01$             & $0.009\pm0.002 $         & $2.266\pm0.004 $        & $2.275\pm0.006$           \\
        NeuralIF & $100.0\pm0.0$         & $9.00\pm0.01$              & $0.732\pm0.005 $         & $1.947\pm0.064 $        & $2.679\pm0.062$           \\
        MPSSMPCG & $100.0\pm0.0$         & $9.23\pm0.02$              & $0.299\pm0.004 $         & $2.015\pm0.053 $        & $2.314\pm0.052$           \\
    \end{tabular}
\end{table}

\begin{table}[H]
    \caption{
        Preconditioning performence on peptides-struct graph systems with
        2000/500/500 split. Mehods are evaluated on their PCG performence on
        unseen test system. Results are averaged over three seeds $\{0,1,2\}$.
    }

    \centering
    \small
    \label{tbl:peptides_struct_preconditioning_err}
    \begin{tabular}{c c c}
        Method   & Rel. Error  $\downarrow$                  & Rel. Reisdual  $\downarrow$               \\
        \hline
        Identity & $2.220\cdot 10^{-6}\pm5.722\cdot 10^{-9}$ & $7.815\cdot 10^{-7}\pm1.628\cdot 10^{-9}$ \\
        Jacobi   & $2.031\cdot 10^{-6}\pm3.599\cdot 10^{-9}$ & $6.932\cdot 10^{-7}\pm1.051\cdot 10^{-9}$ \\
        NeuralIF & $1.679\cdot 10^{-6}\pm8.283\cdot 10^{-9}$ & $6.393\cdot 10^{-7}\pm4.363\cdot 10^{-9}$ \\
        MPSSMPCG & $1.602\cdot 10^{-6}\pm2.746\cdot 10^{-8}$ & $5.961\cdot 10^{-7}\pm6.647\cdot 10^{-9}$ \\
    \end{tabular}
\end{table}

